% Generado con se-ramify. Para incluirlo en vuestra plantilla: \input{...}
% Necesita graphicx, longtable, booktabs, float, xcolor y fvextra.
% se-ramify: an image at its natural size, shrunk to fit the text block when bigger.
\providecommand{\seramifyfigura}[1]{%
  \sbox0{\includegraphics{#1}}%
  \ifdim\wd0>\linewidth
    \includegraphics[width=\linewidth,height=0.78\textheight,keepaspectratio]{#1}%
  \else\ifdim\ht0>0.78\textheight
    \includegraphics[height=0.78\textheight,keepaspectratio]{#1}%
  \else\usebox0\fi\fi}

\section*{Ejemplo — BiblioYa — Entrega 1 — Empresa y TCO}
\addcontentsline{toc}{section}{Ejemplo — BiblioYa — Entrega 1 — Empresa y TCO}

\emph{Exportado con se-ramify el 23/10/2026.}

\subsection*{Contenido de la entrega}

{\small
\begin{longtable}{@{}>{\raggedright\arraybackslash}p{0.437\linewidth}>{\raggedright\arraybackslash}p{0.061\linewidth}>{\raggedright\arraybackslash}p{0.425\linewidth}@{}}
\toprule
\textbf{Apartado} & \textbf{Peso} & \textbf{Diagramas} \\
\midrule
\endhead
Definición de la empresa: nombre, marca, modelo de negocio y roles & 25 \% & E1 · Modelo de negocio — Faro Software; E1 · Organigrama y roles; \emph{Nombre y marca (logo, lema) (texto)} \\
Descripción completa de las dos alternativas, con su arquitectura & 25 \% & E1 · Arquitectura — local; E1 · Arquitectura — nube \\
TCO a 3 años con costes completos y realistas de cada alternativa & 25 \% & E1 · TCO a 3 años \\
Recomendación razonada de la alternativa más adecuada & 25 \% & E1 · TCO a 3 años \\
\bottomrule
\end{longtable}}

\subsection*{1. E1 · Modelo de negocio — Faro Software}

\emph{Lienzo de modelo de negocio}

\begin{figure}[H]
\centering
\seramifyfigura{01-e1-modelo-de-negocio-faro-software.png}
\caption{E1 · Modelo de negocio — Faro Software}
\end{figure}

Faro Software es una empresa inventada de cuatro personas. Este lienzo describe la empresa, no el producto.

\textbf{Negocio:} Faro Software — software a medida para instituciones educativas y culturales

\subsubsection*{Socios clave}

\begin{itemize}
\item Proveedor de nube (alojamiento)
\item Universidades con las que ya trabajamos
\item Freelance de diseño UX
\end{itemize}

\subsubsection*{Actividades clave}

\begin{itemize}
\item Análisis de requisitos con el cliente
\item Desarrollo y pruebas
\item Mantenimiento y soporte
\end{itemize}

\subsubsection*{Recursos clave}

\begin{itemize}
\item Equipo de 4 ingenieros
\item Plantillas y componentes reutilizables
\item Experiencia en el sector educativo
\end{itemize}

\subsubsection*{Propuesta de valor}

\begin{itemize}
\item Aplicaciones a medida entregadas por iteraciones cortas
\item El cliente ve una versión que funciona cada 3 semanas
\item Documentación formal para poder mantenerla
\end{itemize}

\subsubsection*{Relación con clientes}

\begin{itemize}
\item Reunión con el cliente al final de cada iteración
\item Un único interlocutor (jefa de proyecto)
\end{itemize}

\subsubsection*{Canales}

\begin{itemize}
\item Recomendación entre universidades
\item Concursos públicos
\item Web y LinkedIn
\end{itemize}

\subsubsection*{Segmentos de clientes}

\begin{itemize}
\item Bibliotecas y servicios universitarios
\item Museos y archivos pequeños
\end{itemize}

\subsubsection*{Estructura de costes}

\begin{itemize}
\item Salarios ($\approx$80 \%)
\item Licencias y nube para desarrollo
\item Formación
\end{itemize}

\subsubsection*{Fuentes de ingresos}

\begin{itemize}
\item Proyecto a precio cerrado por fases
\item Contrato anual de mantenimiento (15 \% del proyecto)
\item Bolsa de horas para cambios
\end{itemize}

\subsubsection*{Notas}

Faro Software es una empresa inventada de cuatro personas. Este lienzo describe la empresa, no el producto.

\subsection*{2. E1 · Organigrama y roles}

\emph{Organigrama y roles}

\begin{figure}[H]
\centering
\seramifyfigura{02-e1-organigrama-y-roles.png}
\caption{E1 · Organigrama y roles}
\end{figure}

{\footnotesize\setlength{\tabcolsep}{4pt}
\begin{longtable}{@{}>{\raggedright\arraybackslash}p{0.212\linewidth}>{\raggedright\arraybackslash}p{0.106\linewidth}>{\raggedright\arraybackslash}p{0.142\linewidth}>{\raggedright\arraybackslash}p{0.191\linewidth}>{\raggedright\arraybackslash}p{0.263\linewidth}@{}}
\toprule
\textbf{Rol} & \textbf{Persona} & \textbf{Área} & \textbf{Depende de} & \textbf{Responsabilidades} \\
\midrule
\endhead
Dirección y jefatura de proyecto & Ana López &  & — & Planifica y sigue el proyecto (Gantt)\newline{}Interlocutora con el cliente\newline{}Redacta las actas \\
Análisis y requisitos & Bruno Díaz & Equipo de desarrollo & Dirección y jefatura de proyecto & Casos de uso y especificaciones\newline{}Historias de usuario y Gherkin \\
Arquitectura y diseño & Carla Gómez & Equipo de desarrollo & Dirección y jefatura de proyecto & Diagramas de clases y secuencia\newline{}Arquitectura y TCO \\
Desarrollo y calidad & David Martín & Equipo de desarrollo & Dirección y jefatura de proyecto & Implementación\newline{}Plan y casos de prueba \\
\bottomrule
\end{longtable}}

\subsection*{3. E1 · Arquitectura — local}

\emph{Diagrama de arquitectura / equipos}

\begin{figure}[H]
\centering
\seramifyfigura{03-e1-arquitectura-local.png}
\caption{E1 · Arquitectura — local}
\end{figure}

Alternativa A. Todo duplicado donde una avería pararía el servicio: dos servidores con réplica, dos switches, SAI y copias en NAS.

\subsubsection*{Sala técnica de la biblioteca}

{\small
\begin{longtable}{@{}>{\raggedright\arraybackslash}p{0.596\linewidth}>{\raggedright\arraybackslash}p{0.353\linewidth}@{}}
\toprule
\textbf{Equipo} & \textbf{Tipo} \\
\midrule
\endhead
Router + firewall & Cortafuegos \\
SAI & SAI \\
Switch principal ×2 & Switch \\
Servidor 1 (app + BD) & Servidor \\
Servidor 2 (réplica) & Servidor \\
NAS de copias & NAS / copias \\
\bottomrule
\end{longtable}}

\subsubsection*{Mostradores y despachos}

{\small
\begin{longtable}{@{}>{\raggedright\arraybackslash}p{0.534\linewidth}>{\raggedright\arraybackslash}p{0.415\linewidth}@{}}
\toprule
\textbf{Equipo} & \textbf{Tipo} \\
\midrule
\endhead
Switch de planta & Switch \\
Lectores de carné ×4 & TPV \\
PC mostrador ×4 & PC de sobremesa \\
Portátiles gestión ×2 & Portátil \\
\bottomrule
\end{longtable}}

\subsubsection*{Fuera de zonas}

{\small
\begin{longtable}{@{}>{\raggedright\arraybackslash}p{0.686\linewidth}>{\raggedright\arraybackslash}p{0.263\linewidth}@{}}
\toprule
\textbf{Equipo} & \textbf{Tipo} \\
\midrule
\endhead
Socios (móvil) & Móvil \\
\bottomrule
\end{longtable}}

\subsubsection*{Conexiones}

\begin{itemize}
\item Router + firewall $\rightarrow$ Switch principal ×2
\item SAI $\rightarrow$ Servidor 1 (app + BD) (corriente)
\item Switch principal ×2 $\rightarrow$ Servidor 1 (app + BD)
\item Switch principal ×2 $\rightarrow$ Servidor 2 (réplica)
\item Switch principal ×2 $\rightarrow$ NAS de copias
\item Servidor 1 (app + BD) $\rightarrow$ Servidor 2 (réplica) (réplica)
\item Switch principal ×2 $\rightarrow$ Switch de planta
\item Switch de planta $\rightarrow$ PC mostrador ×4
\item Switch de planta $\rightarrow$ Portátiles gestión ×2
\item Lectores de carné ×4 $\rightarrow$ PC mostrador ×4 (USB)
\item Socios (móvil) $\rightarrow$ Router + firewall (Internet (HTTPS))
\end{itemize}

\subsection*{4. E1 · Arquitectura — nube}

\emph{Diagrama de arquitectura / equipos}

\begin{figure}[H]
\centering
\seramifyfigura{04-e1-arquitectura-nube.png}
\caption{E1 · Arquitectura — nube}
\end{figure}

Alternativa B. La aplicación en dos zonas de disponibilidad detrás de un balanceador, con base de datos gestionada multi-AZ.

\subsubsection*{Biblioteca}

{\small
\begin{longtable}{@{}>{\raggedright\arraybackslash}p{0.534\linewidth}>{\raggedright\arraybackslash}p{0.415\linewidth}@{}}
\toprule
\textbf{Equipo} & \textbf{Tipo} \\
\midrule
\endhead
Router + firewall & Cortafuegos \\
Switch & Switch \\
PC mostrador ×4 & PC de sobremesa \\
Portátiles gestión ×2 & Portátil \\
\bottomrule
\end{longtable}}

\subsubsection*{Nube (región eu-south-2, Madrid)}

{\small
\begin{longtable}{@{}>{\raggedright\arraybackslash}p{0.551\linewidth}>{\raggedright\arraybackslash}p{0.398\linewidth}@{}}
\toprule
\textbf{Equipo} & \textbf{Tipo} \\
\midrule
\endhead
Balanceador de carga & Switch \\
VM app — zona A & Servicio en la nube \\
VM app — zona B & Servicio en la nube \\
BD gestionada (multi-AZ) & Base de datos \\
Copias (almacenamiento) & NAS / copias \\
\bottomrule
\end{longtable}}

\subsubsection*{Fuera de zonas}

{\small
\begin{longtable}{@{}>{\raggedright\arraybackslash}p{0.686\linewidth}>{\raggedright\arraybackslash}p{0.263\linewidth}@{}}
\toprule
\textbf{Equipo} & \textbf{Tipo} \\
\midrule
\endhead
Socios (móvil) & Móvil \\
\bottomrule
\end{longtable}}

\subsubsection*{Conexiones}

\begin{itemize}
\item Router + firewall $\rightarrow$ Balanceador de carga (Internet (HTTPS))
\item Switch $\rightarrow$ Router + firewall
\item PC mostrador ×4 $\rightarrow$ Switch
\item Portátiles gestión ×2 $\rightarrow$ Switch
\item Balanceador de carga $\rightarrow$ VM app — zona A
\item Balanceador de carga $\rightarrow$ VM app — zona B
\item VM app — zona A $\rightarrow$ BD gestionada (multi-AZ)
\item VM app — zona B $\rightarrow$ BD gestionada (multi-AZ)
\item BD gestionada (multi-AZ) $\rightarrow$ Copias (almacenamiento) (copia diaria)
\item Socios (móvil) $\rightarrow$ Balanceador de carga (4G / Wi-Fi)
\end{itemize}

\subsection*{5. E1 · TCO a 3 años}

\emph{TCO (coste total de propiedad)}

\begin{figure}[H]
\centering
\seramifyfigura{05-e1-tco-a-3-anos.png}
\caption{E1 · TCO a 3 años}
\end{figure}

\emph{Coste total de propiedad a 3 años, sin inflación ni tasa de descuento. Importes en euros.}

\subsubsection*{Local (on-premises)}

Dos servidores en la sala técnica de la biblioteca central: uno principal y otro de réplica en espera (la base de datos se replica en tiempo real). SAI para cortes de luz, NAS con copias diarias y una copia semanal fuera del edificio. Un cortafuegos y un switch gestionable conectan con los puestos de las cinco bibliotecas.

{\scriptsize\setlength{\tabcolsep}{3pt}
\begin{longtable}{@{}>{\raggedright\arraybackslash}p{0.109\linewidth}>{\raggedright\arraybackslash}p{0.195\linewidth}>{\raggedright\arraybackslash}p{0.071\linewidth}>{\raggedright\arraybackslash}p{0.071\linewidth}>{\raggedright\arraybackslash}p{0.078\linewidth}>{\raggedright\arraybackslash}p{0.069\linewidth}>{\raggedright\arraybackslash}p{0.078\linewidth}>{\raggedright\arraybackslash}p{0.078\linewidth}>{\raggedright\arraybackslash}p{0.136\linewidth}@{}}
\toprule
\textbf{Bloque} & \textbf{Concepto} & \textbf{Tipo} & \textbf{Pago} & \textbf{Año 1} & \textbf{Año 2} & \textbf{Año 3} & \textbf{Total} & \textbf{Fuente} \\
\midrule
\endhead
Desarrollo & Análisis, desarrollo y pruebas (horas del equipo, 35 €/h) (×1200) & CAPEX & Una vez (año 1) & 42.000~€ & 0~€ & 0~€ & \textbf{42.000~€} & ejemplo — horas estimadas × coste/hora \\
Desarrollo & Instalación, configuración y migración de datos & CAPEX & Una vez (año 1) & 3.200~€ & 0~€ & 0~€ & \textbf{3.200~€} & ejemplo — sustituir por precio real \\
Desarrollo & Formación del personal de las bibliotecas & CAPEX & Una vez (año 1) & 1.200~€ & 0~€ & 0~€ & \textbf{1.200~€} & ejemplo — sustituir por precio real \\
HW y SW & Servidor Dell PowerEdge T150 (Xeon E-2314, 32 GB, 2× SSD 960 GB) (×2) & CAPEX & Una vez (año 1) & 3.300~€ & 0~€ & 0~€ & \textbf{3.300~€} & ejemplo — sustituir por precio real \\
HW y SW & SAI APC Smart-UPS 1500 VA & CAPEX & Una vez (año 1) & 780~€ & 0~€ & 0~€ & \textbf{780~€} & ejemplo — sustituir por precio real \\
HW y SW & NAS Synology DS923+ & CAPEX & Una vez (año 1) & 620~€ & 0~€ & 0~€ & \textbf{620~€} & ejemplo — sustituir por precio real \\
HW y SW & Discos para el NAS (4 TB) (×4) & CAPEX & Una vez (año 1) & 480~€ & 0~€ & 0~€ & \textbf{480~€} & ejemplo — sustituir por precio real \\
HW y SW & Switch gestionable 24 puertos & CAPEX & Una vez (año 1) & 280~€ & 0~€ & 0~€ & \textbf{280~€} & ejemplo — sustituir por precio real \\
HW y SW & Cortafuegos FortiGate 40F & CAPEX & Una vez (año 1) & 650~€ & 0~€ & 0~€ & \textbf{650~€} & ejemplo — sustituir por precio real \\
HW y SW & Armario rack 12U & CAPEX & Una vez (año 1) & 350~€ & 0~€ & 0~€ & \textbf{350~€} & ejemplo — sustituir por precio real \\
HW y SW & Windows Server 2025 Standard (16 núcleos) (×2) & CAPEX & Una vez (año 1) & 2.300~€ & 0~€ & 0~€ & \textbf{2.300~€} & ejemplo — sustituir por precio real \\
HW y SW & Ampliación de garantía del fabricante (×2) & OPEX & Al año desde el año 2 & 0~€ & 380~€ & 380~€ & \textbf{760~€} & ejemplo — sustituir por precio real \\
Operaciones & Administrador de sistemas (20 \% de 38.000 €/año) & OPEX & Al año & 7.600~€ & 7.600~€ & 7.600~€ & \textbf{22.800~€} & ejemplo — tabla salarial \\
Operaciones & Copia externa en la nube (2 TB) & OPEX & Al mes & 300~€ & 300~€ & 300~€ & \textbf{900~€} & ejemplo — sustituir por precio real \\
Operaciones & Mantenimiento correctivo y parches (40 h/año) & OPEX & Al año & 1.400~€ & 1.400~€ & 1.400~€ & \textbf{4.200~€} & ejemplo — sustituir por precio real \\
Largo plazo & Reemplazo de discos del NAS (×4) & CAPEX & Una vez (año 3) & 0~€ & 0~€ & 480~€ & \textbf{480~€} & ejemplo — sustituir por precio real \\
 & \textbf{Total} &  &  & \textbf{64.460~€} & \textbf{9.680~€} & \textbf{10.160~€} & \textbf{84.300~€} &  \\
\bottomrule
\end{longtable}}

\subsubsection*{Nube (AWS, región España)}

La aplicación corre en dos máquinas virtuales en dos zonas de disponibilidad detrás de un balanceador; la base de datos es PostgreSQL gestionado con réplica en otra zona. Copias automáticas diarias. Las bibliotecas se conectan por Internet (HTTPS) y la sede, además, por VPN.

{\scriptsize\setlength{\tabcolsep}{3pt}
\begin{longtable}{@{}>{\raggedright\arraybackslash}p{0.109\linewidth}>{\raggedright\arraybackslash}p{0.201\linewidth}>{\raggedright\arraybackslash}p{0.071\linewidth}>{\raggedright\arraybackslash}p{0.072\linewidth}>{\raggedright\arraybackslash}p{0.078\linewidth}>{\raggedright\arraybackslash}p{0.069\linewidth}>{\raggedright\arraybackslash}p{0.069\linewidth}>{\raggedright\arraybackslash}p{0.078\linewidth}>{\raggedright\arraybackslash}p{0.138\linewidth}@{}}
\toprule
\textbf{Bloque} & \textbf{Concepto} & \textbf{Tipo} & \textbf{Pago} & \textbf{Año 1} & \textbf{Año 2} & \textbf{Año 3} & \textbf{Total} & \textbf{Fuente} \\
\midrule
\endhead
Desarrollo & Análisis, desarrollo y pruebas (horas del equipo, 35 €/h) (×1200) & CAPEX & Una vez (año 1) & 42.000~€ & 0~€ & 0~€ & \textbf{42.000~€} & ejemplo — horas estimadas × coste/hora \\
Desarrollo & Migración de datos y puesta en marcha en la nube & CAPEX & Una vez (año 1) & 2.400~€ & 0~€ & 0~€ & \textbf{2.400~€} & ejemplo — sustituir por precio real \\
Desarrollo & Formación del personal de las bibliotecas & CAPEX & Una vez (año 1) & 1.200~€ & 0~€ & 0~€ & \textbf{1.200~€} & ejemplo — sustituir por precio real \\
HW y SW & Máquinas virtuales EC2 t3.medium (una por zona) (×2) & OPEX & Al mes & 840~€ & 840~€ & 840~€ & \textbf{2.520~€} & ejemplo — calculadora de AWS \\
HW y SW & Base de datos RDS PostgreSQL db.t3.medium multi-AZ & OPEX & Al mes & 1.680~€ & 1.680~€ & 1.680~€ & \textbf{5.040~€} & ejemplo — calculadora de AWS \\
HW y SW & Balanceador de carga (ALB) & OPEX & Al mes & 264~€ & 264~€ & 264~€ & \textbf{792~€} & ejemplo — calculadora de AWS \\
HW y SW & Almacenamiento de copias (S3, 300 GB) & OPEX & Al mes & 96~€ & 96~€ & 96~€ & \textbf{288~€} & ejemplo — calculadora de AWS \\
HW y SW & Tráfico de salida & OPEX & Al mes & 60~€ & 60~€ & 60~€ & \textbf{180~€} & ejemplo — calculadora de AWS \\
HW y SW & VPN sitio a sitio con la biblioteca central & OPEX & Al mes & 432~€ & 432~€ & 432~€ & \textbf{1.296~€} & ejemplo — calculadora de AWS \\
Operaciones & Plan de soporte Business & OPEX & Al mes & 1.200~€ & 1.200~€ & 1.200~€ & \textbf{3.600~€} & ejemplo — web del proveedor \\
Operaciones & Administrador cloud (10 \% de 40.000 €/año) & OPEX & Al año & 4.000~€ & 4.000~€ & 4.000~€ & \textbf{12.000~€} & ejemplo — tabla salarial \\
Operaciones & Monitorización y alertas & OPEX & Al mes & 120~€ & 120~€ & 120~€ & \textbf{360~€} & ejemplo — calculadora de AWS \\
Largo plazo & Crecimiento de uso (más socios y préstamos) & OPEX & Al año desde el año 2 & 0~€ & 300~€ & 300~€ & \textbf{600~€} & ejemplo — sustituir por precio real \\
Largo plazo & Plan de salida: exportar datos si se cambia de proveedor & CAPEX & Una vez (año 3) & 0~€ & 0~€ & 800~€ & \textbf{800~€} & ejemplo — sustituir por precio real \\
 & \textbf{Total} &  &  & \textbf{54.292~€} & \textbf{8.992~€} & \textbf{9.792~€} & \textbf{73.076~€} &  \\
\bottomrule
\end{longtable}}

\subsubsection*{Resumen}

{\scriptsize\setlength{\tabcolsep}{3pt}
\begin{longtable}{@{}>{\raggedright\arraybackslash}p{0.253\linewidth}>{\raggedright\arraybackslash}p{0.083\linewidth}>{\raggedright\arraybackslash}p{0.073\linewidth}>{\raggedright\arraybackslash}p{0.083\linewidth}>{\raggedright\arraybackslash}p{0.123\linewidth}>{\raggedright\arraybackslash}p{0.083\linewidth}>{\raggedright\arraybackslash}p{0.083\linewidth}>{\raggedright\arraybackslash}p{0.116\linewidth}@{}}
\toprule
\textbf{Alternativa} & \textbf{Año 1} & \textbf{Año 2} & \textbf{Año 3} & \textbf{Total 3 años} & \textbf{CAPEX} & \textbf{OPEX} & \textbf{Por usuario} \\
\midrule
\endhead
Local (on-premises) & 64.460~€ & 9.680~€ & 10.160~€ & \textbf{84.300~€} & 55.640~€ & 28.660~€ & 2.810,00~€ \\
Nube (AWS, región España) & 54.292~€ & 8.992~€ & 9.792~€ & \textbf{73.076~€} & 46.400~€ & 26.676~€ & 2.435,87~€ \\
\bottomrule
\end{longtable}}

{\small
\begin{longtable}{@{}>{\raggedright\arraybackslash}p{0.286\linewidth}>{\raggedright\arraybackslash}p{0.268\linewidth}>{\raggedright\arraybackslash}p{0.369\linewidth}@{}}
\toprule
\textbf{Bloque} & \textbf{Local (on-premises)} & \textbf{Nube (AWS, región España)} \\
\midrule
\endhead
Desarrollo y despliegue & 46.400~€ & 45.600~€ \\
Hardware y software & 9.520~€ & 10.116~€ \\
Operaciones & 27.900~€ & 15.960~€ \\
Gastos a largo plazo & 480~€ & 1.400~€ \\
\bottomrule
\end{longtable}}

La alternativa más barata a 3 años es \textbf{Nube (AWS, región España)} (73.076~€), 11.224~€ menos que Local (on-premises).

Calculado para 30 usuarios.

\subsubsection*{Comparación cualitativa}

\emph{Puntuación de 1 (peor) a 5 (mejor), ponderada por el peso de cada criterio.}

{\small
\begin{longtable}{@{}>{\raggedright\arraybackslash}p{0.297\linewidth}>{\raggedright\arraybackslash}p{0.061\linewidth}>{\raggedright\arraybackslash}p{0.230\linewidth}>{\raggedright\arraybackslash}p{0.310\linewidth}@{}}
\toprule
\textbf{Criterio} & \textbf{Peso} & \textbf{Local (on-premises)} & \textbf{Nube (AWS, región España)} \\
\midrule
\endhead
Disponibilidad (9–18 h, L–V) & 30 \% & 3 & 5 \\
Redundancia & 30 \% & 3 & 5 \\
Riesgo operativo & 20 \% & 2 & 4 \\
Escalabilidad & 10 \% & 2 & 5 \\
Dependencia del proveedor & 10 \% & 5 & 2 \\
\textbf{Puntuación ponderada} &  & \textbf{2,90} & \textbf{4,50} \\
\bottomrule
\end{longtable}}

\subsubsection*{Recomendación}

Recomendamos la alternativa en la nube. A tres años sale unos 11.000 € más barata, sobre todo porque no hay que comprar servidores ni dedicar tanto tiempo de administración, y cumple mejor los requisitos de la biblioteca: si cae una máquina o una zona, la otra sigue atendiendo, y el proveedor garantiza la disponibilidad en horario de apertura sin depender de una sala técnica que se calienta. Su punto débil es la dependencia del proveedor; por eso incluimos un plan de salida (exportar los datos en formato abierto) y usamos PostgreSQL, que se puede llevar a otro proveedor o a local. Si la biblioteca no tuviera buena conexión a Internet, la alternativa local volvería a ser razonable.
